\newpage
\chapter{METODOLOGÍA}

% ------------------------------------------------------------
\section{Metodología}
% ------------------------------------------------------------

El presente proyecto adopta \textbf{Scrum} como metodología de desarrollo, un marco de
trabajo ágil ampliamente utilizado en proyectos de software que demandan flexibilidad,
entregas incrementales y colaboración continua entre los integrantes del equipo.

\subsection{Justificación de la elección}

Scrum resulta idóneo para este proyecto por las siguientes razones:

\begin{itemize}
    \item \textbf{Adaptabilidad a cambios:} los requerimientos del TUPA pueden evolucionar
            durante el levantamiento de información con las oficinas de la UNSAAC; Scrum
            permite incorporar estos cambios de forma ordenada en cada sprint.
    \item \textbf{Entregas rápidas e incrementales:} al dividir el desarrollo en sprints de
            duración fija, se obtienen avances funcionales verificables al final de cada
            iteración, facilitando la retroalimentación oportuna del docente y los usuarios.
    \item \textbf{Trabajo colaborativo:} el equipo está conformado por cinco integrantes con
            roles complementarios; Scrum define responsabilidades claras y promueve la
            comunicación constante mediante ceremonias estructuradas.
    \item \textbf{Transparencia y seguimiento:} el uso de tableros de tareas y herramientas
            como Jira permite visualizar el estado del proyecto en todo momento, favoreciendo
            la trazabilidad de actividades y el cumplimiento de los indicadores establecidos
            en el plan de monitoreo.
\end{itemize}

\subsection{Fases del ciclo de desarrollo}

El ciclo de vida del proyecto se estructura en las siguientes fases, ejecutadas de forma
iterativa a lo largo de los sprints planificados:

\begin{enumerate}
    \item \textbf{Planificación:} definición del alcance del sprint, estimación de historias
            de usuario y asignación de tareas al equipo.
    \item \textbf{Análisis:} levantamiento y refinamiento de requerimientos funcionales y no
            funcionales del sistema, en coordinación con las oficinas administrativas de la
            UNSAAC.
    \item \textbf{Diseño:} elaboración de prototipos de interfaz en Figma, diseño de la
            arquitectura del sistema y modelado de la base de datos.
    \item \textbf{Implementación:} desarrollo de los módulos frontend y backend conforme a
            las historias de usuario aprobadas en cada sprint.
    \item \textbf{Pruebas:} verificación funcional, pruebas de integración y validación de
            usabilidad con usuarios representativos.
    \item \textbf{Despliegue:} publicación de incrementos en el entorno de producción y
            monitoreo de disponibilidad y rendimiento del sistema.
\end{enumerate}

\subsection{Herramientas de gestión}

El seguimiento del proyecto se realizará principalmente mediante \textbf{Jira}, herramienta
que permite gestionar el backlog del producto, planificar sprints, registrar el avance de
tareas y controlar los indicadores de desempeño definidos en el plan de monitoreo. De
forma complementaria, se utilizará \textbf{GitHub} para el control de versiones del código
fuente y la gestión de incidencias técnicas.

\subsection{Roles del equipo Scrum}

La siguiente tabla describe los roles Scrum asignados a cada participante del proyecto:

\begin{table}[H]
\centering
\caption{Roles Scrum del equipo de desarrollo}
\label{tab:roles_scrum}
\begin{tabular}{|p{5cm}|p{3.5cm}|p{6cm}|}
\hline
\textbf{Integrante} & \textbf{Rol Scrum} & \textbf{Responsabilidad principal} \\
\hline
Ing. Julio Vladimir Quispe Sota & Product Owner &
Definir y priorizar el backlog del producto; validar que los entregables cumplan con
las necesidades administrativas del TUPA. \\
\hline
Ramos Alata Hernan Washington & Scrum Master &
Facilitar las ceremonias Scrum, eliminar impedimentos y velar por el cumplimiento
del proceso ágil en el equipo. \\
\hline
Vitorino Marin Efrain & Developer – Backend &
Desarrollar la lógica de negocio, API REST y gestión de base de datos; delegar y
coordinar subtareas del módulo backend. \\
\hline
Mamani Flores Natan & Developer – Frontend &
Implementar las interfaces de usuario; delegar y coordinar subtareas del módulo
frontend en coordinación con el diseñador UX/UI. \\
\hline
Cabeza Huillca Flavio Antony & Developer – UX/UI &
Diseñar prototipos en Figma, definir la guía de estilos y coordinar la experiencia
de usuario en toda la plataforma. \\
\hline
Polo Chura Marco Rosauro & QA Tester &
Diseñar y ejecutar casos de prueba, reportar defectos en Jira y verificar el
cumplimiento de los criterios de aceptación por sprint. \\
\hline
\end{tabular}
\end{table}

\subsection{Ceremonias y seguimiento}

El equipo ejecutará las siguientes ceremonias Scrum a lo largo del proyecto:

\begin{itemize}
    \item \textbf{Daily meeting (reunión diaria):} sesión de 15 minutos en la que cada
            integrante comparte lo que realizó el día anterior, lo que hará ese día y los
            impedimentos que enfrenta. Permite detectar bloqueos de forma temprana.
    \item \textbf{Sprint planning (planificación del sprint):} al inicio de cada sprint, el
            equipo selecciona las historias de usuario del backlog, las estima en puntos y las
            descompone en tareas concretas asignadas a cada desarrollador en Jira.
    \item \textbf{Sprint review (revisión del sprint):} al cierre de cada sprint, se
            presenta el incremento funcional al Product Owner y, cuando corresponda, a los
            docentes responsables, para obtener retroalimentación formal.
    \item \textbf{Sprint retrospective (retrospectiva):} el equipo reflexiona sobre el proceso
            de trabajo, identifica áreas de mejora y define acciones concretas para el
            siguiente sprint.
\end{itemize}

% ------------------------------------------------------------
\section{Actividades principales}
% ------------------------------------------------------------

A continuación se describen las actividades principales del proyecto, organizadas conforme
al cronograma establecido. Cada actividad es liderada por un integrante responsable, quien
además delega y coordina las subtareas correspondientes entre los demás miembros del equipo.

\begin{description}

    \item[Actividad 1: Reunión y planificación del proyecto.]
    Consiste en la constitución formal del equipo, la definición del alcance general del
    proyecto, la elaboración del cronograma de actividades y la configuración de las
    herramientas de gestión (Jira, GitHub). Se establecen los acuerdos de comunicación
    interna y se realiza la primera sesión de Sprint Planning. Esta actividad es liderada
    por el \textbf{Scrum Master (Ramos Alata Hernan Washington)}, quien coordina la
    organización del equipo y la asignación inicial de roles y responsabilidades.

    \item[Actividad 2: Levantamiento de requerimientos.]
    Comprende la identificación y documentación de los requerimientos funcionales y no
    funcionales del sistema mediante entrevistas semiestructuradas con las oficinas
    administrativas de la UNSAAC, análisis del TUPA vigente y revisión de documentación
    institucional. Los resultados se registran como historias de usuario en el backlog de
    Jira. Esta actividad es liderada por el \textbf{QA Tester (Polo Chura Marco Rosauro)},
    quien dirige el análisis de requerimientos y delega la elaboración de fichas de
    entrevista y documentos de especificación al resto del equipo.

    \item[Actividad 3: Diseño de prototipos en Figma.]
    Se elaboran los wireframes y prototipos interactivos de alta fidelidad de la plataforma,
    incluyendo la guía de estilos, el sistema de componentes visuales y los flujos de
    navegación para cada tipo de usuario (estudiante, personal administrativo y jefe de
    oficina). Los prototipos son validados con usuarios representativos antes de iniciar el
    desarrollo. Esta actividad es liderada por el \textbf{desarrollador UX/UI (Cabeza
    Huillca Flavio Antony)}, quien diseña los prototipos y coordina los ajustes con el
    equipo frontend.

    \item[Actividad 4: Diseño de la base de datos.]
    Incluye el modelado conceptual, lógico y físico de la base de datos del sistema,
    definiendo entidades, relaciones, restricciones de integridad y procedimientos
    almacenados necesarios para soportar los módulos de gestión de trámites, usuarios y
    reportes. El modelo es revisado y aprobado por el Product Owner antes de su
    implementación. Esta actividad es liderada por el \textbf{desarrollador backend
    (Vitorino Marin Efrain)}, quien diseña el esquema de la base de datos y delega la
    documentación del modelo al equipo.

    \item[Actividad 5: Desarrollo de interfaces gráficas e implementación del módulo
    frontend.]
    Comprende la traducción de los prototipos aprobados en Figma a código, implementando
    las vistas, componentes y flujos de interacción de la plataforma web. Se aplican
    buenas prácticas de desarrollo frontend, diseño responsivo y accesibilidad. Esta
    actividad es liderada por el \textbf{desarrollador frontend (Mamani Flores Natan)},
    quien coordina la implementación de interfaces y delega el desarrollo de componentes
    específicos a los demás integrantes.

    \item[Actividad 6: Integración de módulos e implementación de notificaciones y
    validaciones.]
    Se integran los módulos frontend y backend del sistema, incluyendo el motor de
    notificaciones (correo electrónico y notificaciones en plataforma) y los mecanismos
    de validación documental. Se verifican los contratos de la API REST y la correcta
    comunicación entre componentes. Esta actividad es liderada por el \textbf{desarrollador
    backend (Vitorino Marin Efrain)}, quien coordina la integración técnica y gestiona los
    ajustes necesarios con el equipo frontend.

    \item[Actividad 7: Pruebas y corrección de errores.]
    Abarca la ejecución de pruebas funcionales, de integración y de usabilidad sobre los
    módulos desarrollados. Los defectos identificados se registran en Jira y se priorizan
    para su corrección en el sprint correspondiente. Se verifica el cumplimiento de los
    criterios de aceptación definidos en las historias de usuario. Esta actividad es
    liderada por el \textbf{QA Tester (Polo Chura Marco Rosauro)}, quien diseña los casos
    de prueba y coordina las sesiones de validación con usuarios finales.

    \item[Actividad 8: Elaboración de documentación final y presentación del proyecto.]
    Comprende la redacción de la documentación técnica y académica del proyecto (memoria
    del sistema, manual de usuario, informe final), la preparación de la presentación ante
    los docentes y la entrega formal de todos los artefactos del proyecto. Esta actividad
    es liderada por el \textbf{Scrum Master (Ramos Alata Hernan Washington)}, quien
    coordina la consolidación de la documentación y delega la redacción de secciones
    específicas a cada integrante según su área de responsabilidad.

\end{description}